\documentclass[letterpaper]{article}
\usepackage[margin=1in]{geometry}
\usepackage{amsmath,amssymb,amsthm}
\usepackage{hyperref}
\usepackage{booktabs}
\usepackage{graphicx}
\usepackage{listings}
\usepackage{xcolor}
\usepackage{natbib}

\hypersetup{colorlinks=true,linkcolor=blue,citecolor=blue,urlcolor=blue}

\lstset{
  basicstyle=\ttfamily\small,
  breaklines=true,
  frame=single,
  numbers=left,
  numberstyle=\tiny\color{gray},
  xleftmargin=2em,
}

\newtheorem{definition}{Definition}
\newtheorem{theorem}{Theorem}

\title{Book-Guided, LLM-Detailed, Kernel-Verified:\\
Formal Proof Construction from ZFC Axioms to Program Correctness}

\author{Y}
\date{2026}

\begin{document}
\maketitle

\begin{abstract}
Formal proofs are too large for humans to write, read, or verify by inspection.
A two-line definition of addition on natural numbers expands to over 11,500 lines
of formal proof from ZFC axioms. This 1000:1 ratio is not a tooling limitation
--- it is the inherent nature of formal verification.

We present a pipeline for formal proof construction that addresses this gap:
a mathematical textbook provides proof strategy and development order
(\emph{book-guided}); a large language model translates informal reasoning into
formal proof steps (\emph{LLM-detailed}); and a minimal proof kernel guarantees
correctness (\emph{kernel-verified}).

We demonstrate this pipeline by constructing a verified proof chain from ZFC
set theory axioms to the correctness of a Turing machine that computes addition
--- spanning 10 abstraction levels, 203 theorems, and 39,059 lines of proof,
all verified by a 334-line sequent calculus kernel. During formalization, the
kernel caught a definition error in the guiding textbook that three years of
human review had missed.
\end{abstract}

% ============================================================
\section{Introduction}
\label{sec:intro}
% ============================================================

The promise of formal verification is absolute correctness: if a proof checks,
the theorem holds. The barrier is cost. Formal proofs in systems like Coq,
Lean, Isabelle, and Metamath are orders of magnitude longer than their textbook
counterparts, requiring extensive manual effort from trained experts.

Large language models (LLMs) offer a potential solution. They can generate
large volumes of plausible formal text, and they can work at a scale of
mechanical detail that humans cannot sustain. But LLMs are unreliable ---
they hallucinate, make logical errors, and produce subtly wrong proofs. Used
alone, they cannot be trusted for formal verification.

We propose a pipeline that combines three components:

\begin{enumerate}
  \item \textbf{Book-guided.} A mathematical textbook, written with formal
    modeling in mind, provides the proof strategy, definitions, and
    development order. The book is the human-readable specification of what
    should be proved and why.
  \item \textbf{LLM-detailed.} A large language model translates the
    textbook's informal reasoning into formal proof steps. It fills in the
    mechanical detail, proposes intermediate lemmas when the proof gets stuck,
    and navigates the growing theorem library by semantic association.
  \item \textbf{Kernel-verified.} A minimal proof kernel checks every proof
    step at construction time. If the kernel accepts a proof, it is correct
    regardless of whether the LLM made errors along the way. The kernel is
    small enough (334 lines) to be audited by a human.
\end{enumerate}

The key insight is that these three components address different aspects of a
structural problem. Formal proofs are humanly impossible to write (too long),
humanly impossible to read (too detailed), and humanly impossible to verify by
inspection (too complex). The textbook is the human interface. The LLM is the
translator. The kernel is the guarantor.

We demonstrate this pipeline by constructing a complete formal proof that a
specific Turing machine correctly computes addition in unary representation,
starting from the axioms of Zermelo set theory (Z).
The proof chain spans 10 abstraction levels --- from pure logic through set
theory, ordered pairs, functions, natural numbers, recursion, arithmetic,
Turing machine semantics, tape operations, and program correctness. It
comprises 203 theorems in 39,059 lines of Python, all verified by a 334-line
sequent calculus kernel.

During the formalization process, the kernel caught a definition error in the
guiding textbook (written by the first author three years prior): the
recursion theorem's definition of ``recursive function'' lacked a domain
constraint, making uniqueness unprovable. Three years of human review had
missed this. The kernel rejected the proof mechanically.

% ============================================================
\section{The Scale of Formal Proof}
\label{sec:scale}
% ============================================================

The gap between informal mathematics and formal proof is not a small constant
factor. To illustrate, consider the definition of addition on natural numbers
in ZFC set theory.

\paragraph{Textbook definition.} Two lines:
\begin{align*}
  m + 0 &= m \\
  m + S(n) &= S(m + n)
\end{align*}

\paragraph{Formal proof.} The proof that this function exists, is unique,
commutes, associates, and computes concrete results such as $4+3=7$ occupies
11,532 lines of formal proof code. This covers a single arithmetic operation
--- addition --- and does not include the prerequisite infrastructure (ordered
pairs, functions, natural numbers, the recursion theorem) which adds
approximately 28,000 more lines.

The ratio is roughly 1000:1. A page of textbook mathematics becomes a
thousand lines of formal proof. This ratio is not specific to our system; it
reflects the inherent gap between informal mathematical reasoning and the
level of detail required for machine verification. Every implicit appeal to
``it is easy to see,'' every omitted case analysis, every use of a previously
established fact must be made explicit.

At this scale, three human limitations become absolute:

\begin{itemize}
  \item \textbf{Humans cannot write these proofs.} The sheer volume of
    mechanical proof steps exceeds what a person can sustain. Writing 39,059
    lines of correct proof code by hand, where each line must satisfy a proof
    kernel's rule checking, is not difficult --- it is infeasible.
  \item \textbf{Humans cannot read these proofs.} The proof of a single
    arithmetic lemma may span hundreds of lines of tactic applications,
    variable instantiations, and cut compositions. No human can hold this in
    working memory or extract meaning from it by inspection.
  \item \textbf{Humans cannot verify these proofs.} Without a mechanical
    checker, there is no way to confirm that 39,059 lines of proof are
    correct. The proof exists beyond human verification capacity.
\end{itemize}

These are not limitations of proof engineering or tooling. They are structural
consequences of what formal proof \emph{is}. Any system that aspires to
large-scale formal verification must reckon with this gap.

To make this concrete: the complete proof of TM addition correctness requires
134,089 individual kernel rule applications (counting each shared lemma once).
These span 150 unique theorems across 10 abstraction levels. Each rule
application is a sequent calculus step --- an axiom introduction, a
quantifier instantiation, a cut, a weakening --- mechanically validated by
the kernel at construction time. No human could perform or check 134,089
sequential logical steps by hand.

% ============================================================
\section{Architecture}
\label{sec:arch}
% ============================================================

\subsection{Foundation: ZFC and Sequent Calculus}

We choose Zermelo-Fraenkel set theory with Choice (ZFC) as the logical
foundation, implemented via Gentzen's sequent calculus (LK).

The choice of ZFC over dependent type theories (the Calculus of Inductive
Constructions used by Coq and Lean) is deliberate. Scientific and
mathematical knowledge is overwhelmingly expressed in the language of sets,
functions, and relations. ZFC is the native foundation of working mathematics.
Encoding mathematical objects in dependent type theory introduces a
translation layer --- natural numbers become inductive types, functions become
dependent products --- that adds complexity without adding expressive power
for the problems we target.

Sequent calculus provides a clean proof-theoretic framework with well-understood
meta-properties (cut elimination, the subformula property). Proofs are
explicit derivation trees, not tactic scripts, making verification
straightforward.

\subsection{The Proof Kernel}

The proof kernel is a 334-line Python module. The complete trusted core
--- formula language (88 lines), ZFC axiom definitions (144 lines),
derived connectives (119 lines), and the proof checker itself ---
totals 685 lines. The kernel implements:

\begin{itemize}
  \item \textbf{Sequent representation.} A sequent $\Gamma \vdash \Delta$
    consists of left (assumptions) and right (conclusions) formula lists with
    set semantics (no duplicates, checked via alpha-equivalence).
  \item \textbf{Eleven proof rules.} Logical rules (axiom, not-left,
    not-right, implies-left, implies-right, forall-left, forall-right, cut)
    and structural rules (weakening-left, weakening-right). Each rule is
    validated at proof construction time.
  \item \textbf{Lazy definition expansion.} Derived connectives (exists, and,
    or, iff, equality) and vocabulary predicates expand to primitive formulas
    on demand.
  \item \textbf{Alpha-equivalence.} Formula comparison respects bound variable
    renaming.
\end{itemize}

Every \texttt{Proof} object validates its rule application at construction.
If the rule is misapplied --- wrong principal formula, eigenvariable violation,
context mismatch --- construction raises an exception. There is no way to
hold an invalid \texttt{Proof} object.

The final \texttt{qed()} function checks that a completed proof has exactly
one formula on the right (the theorem) and only ZFC axiom instances on the
left.

\subsection{Vocabulary Expansion}

Between the kernel and the theorems sits a vocabulary layer: named predicates
that expand to first-order formulas. For example:

\begin{itemize}
  \item $\text{Successor}(s, n)$ expands to
    $\forall x.\, x \in s \leftrightarrow (x \in n \lor x = n)$
  \item $\text{Plus}(m, n, p)$ expands via the recursion theorem's
    characterization of the addition function
  \item $\text{TMConfig}(c, q, h, t)$ expands to an ordered-tuple encoding
    of Turing machine configurations
\end{itemize}

This layer is \emph{outside} the trusted kernel. Vocabulary definitions can
be wrong --- but if they are, the proof will fail to go through. The kernel
treats expanded formulas the same as any other formula.

\subsection{LLM Integration}

The LLM operates as an agent that edits proof source files. It reads theorem
statements and existing proofs, proposes proof steps in Python, and runs the
goal-checking harness. When the kernel rejects a step, the LLM reads the
error message and adjusts.

The LLM's role is \emph{search}: selecting which theorems to apply, in what
order, with what instantiations. The kernel's role is \emph{verification}:
checking that each selected step is logically valid. The LLM can be wrong
arbitrarily often; the kernel catches every error.

In practice, we observed that the LLM navigates the theorem library by
\emph{name semantics} --- guessing that a theorem called
\texttt{plus\_val\_unique} proves uniqueness of addition values, attempting to
use it, and correcting the instantiation when the kernel rejects the first
attempt. This pattern of semantic guessing followed by mechanical error
correction proved effective as the theorem library grew.

% ============================================================
\section{Demonstration: ZFC to Program Correctness}
\label{sec:demo}
% ============================================================

We construct a complete formal proof of the following:

\begin{theorem}
Given a Turing machine $M$ with transition function $\delta$, initial state
$q_0$, and halting state $q_H$: for all natural numbers $a$ and $b$, if $M$
starts with a unary tape encoding of $a$ and $b$ (a ones, separator zero, b
ones), then $M$ halts with a tape encoding of $a + b$.
\end{theorem}

The proof is deliberately restricted to Zermelo set theory (Z):
Extensionality, Empty Set, Pairing, Union, Power Set, Separation, and
Infinity. Replacement, Regularity, and Choice are not used.

\subsection{Abstraction Levels}

The proof chain crosses the following levels, each building on the previous:

\begin{enumerate}
  \item \textbf{Logic.} Sequent calculus rules: modus ponens, double negation,
    quantifier instantiation, if-and-only-if introduction and elimination.
    (24 theorems)
  \item \textbf{Set theory.} Empty set uniqueness, singleton and pair
    existence, union, power set, subset relations.
    (40 theorems)
  \item \textbf{Ordered pairs.} Kuratowski encoding $\langle a,b \rangle =
    \{\{a\},\{a,b\}\}$, injection (if $\langle a,b \rangle = \langle c,d
    \rangle$ then $a=c$ and $b=d$), existence. (included in sets)
  \item \textbf{Functions.} Relations, single-valuedness, domain, range,
    application, function extensionality. (included in sets)
  \item \textbf{Natural numbers.} Omega as the smallest inductive set,
    successor properties, transitive sets, omega induction.
    (9 theorems)
  \item \textbf{Recursion.} The recursion theorem: for any function $f$ and
    initial value $a$, there exists a unique function $h$ with $\text{dom}(h)
    = \omega$, $h(0) = a$, and $h(n^+) = f(h(n))$.
    (36 theorems)
  \item \textbf{Arithmetic.} Addition as a recursive function on $\omega$.
    Existence, uniqueness, commutativity, associativity, concrete computations.
    (44 theorems)
  \item \textbf{TM definitions.} States, transitions, configurations, tapes
    as functions from positions to symbols.
  \item \textbf{Tape operations.} Tape update existence, reading, function
    preservation across updates.
  \item \textbf{Program correctness.} Five execution phases, each proved
    correct by omega induction or single-step verification, composed via a
    trace composition theorem (TMReachesCompose).
    (41 theorems)
\end{enumerate}

\subsection{Summary}

\begin{table}[h]
\centering
\begin{tabular}{lr}
\toprule
\textbf{Metric} & \textbf{Value} \\
\midrule
ZFC axioms defined & 10 \\
Axioms actually used & 7 (Z only) \\
Proof rules in kernel & 11 \\
Proof kernel & 334 lines \\
Trusted core (kernel + language + axioms + derived) & 685 lines \\
Total theorems & 203 \\
Total proof code & 39,059 lines \\
Unique theorems in plus\_comm & 30 \\
Kernel proof steps for plus\_comm (DAG) & 11,593 \\
Unique theorems in tm\_add\_correct & 150 \\
Kernel proof steps for tm\_add\_correct (DAG) & 134,089 \\
Proof-to-kernel ratio & 117:1 \\
Development time & $\sim$25 days \\
Abstraction levels crossed & 10 \\
\bottomrule
\end{tabular}
\caption{Summary statistics of the ZFC-to-TM-correctness proof chain.}
\label{tab:summary}
\end{table}

% ============================================================
\section{Book-Guided Formalization}
\label{sec:book}
% ============================================================

The proof development followed a textbook written by the first author over
three years: a 219-page treatment of mathematical logic and software
engineering covering propositional logic, first-order logic (including sequent
calculus, cut elimination, and completeness), computability, first-order
arithmetic (including G\"{o}del's incompleteness theorems), ZFC set theory
(including ordinals, cardinals, and the axiom of choice), and inner models
(the constructible universe and relative consistency).

The engine's proof kernel implements the sequent calculus from Chapter~3 of
this textbook. The theorem development follows Chapters~4.1 (ZFC axioms,
basic set operations, ordered pairs, functions) and 4.2 (natural numbers,
recursion). The Turing machine correctness proof extends beyond the textbook's
content.

\subsection{The Formalization Caught a Textbook Error}

Theorem~4.2.14 of the textbook defines a recursive function as follows:

\begin{definition}[Textbook version]
$\text{isRecursive}(h, a, f)$ iff $\text{isFunction}(h) \land h(0) = a
\land \forall n \in \omega\; h(n^+) = f(h(n))$.
\end{definition}

This definition contains no domain constraint on $h$. The textbook's proof
derives $\text{dom}(h) = \omega$ as a side result (page~142) but does not
include it in the definition.

When the LLM attempted to formalize the uniqueness proof ($\exists! h\;
\text{isRecursive}(h, a, f)$), the kernel rejected it. The reason:
without $\text{dom}(h) = \omega$, two functions that agree on $\omega$ but
differ on extraneous pairs both satisfy the definition. Uniqueness genuinely
fails under the textbook's formulation.

The fix was to strengthen the definition:

\begin{definition}[Corrected version]
$\text{isRecursive}(h, a, f)$ iff $\text{isFunction}(h) \land
\text{dom}(h) = \omega \land h(0) = a \land \forall n \in \omega\;
h(n^+) = f(h(n))$.
\end{definition}

This error is characteristic of informal mathematics. The textbook proof
works because the reader implicitly restricts attention to the ``intended''
function $h$. The kernel does not permit implicit restrictions. The
definition must be self-contained.

The textbook was written by the first author and reviewed over three years.
The error was caught mechanically, in minutes, by a 334-line kernel during an
LLM-driven formalization attempt.

% ============================================================
\section{Cross-Abstraction Depth}
\label{sec:depth}
% ============================================================

Existing formal verification tools typically operate at a single abstraction
level. SMT solvers (Z3) handle decidable fragments. Interactive theorem
provers (Coq, Lean) work within type theory. Model checkers (TLA+) reason
about state machines. When a problem spans multiple levels --- as most real
systems do --- these tools must be used in isolation, with the cross-level
reasoning left to the human.

The proof of TM addition correctness is a concrete counterexample: a single
proof chain, in a single system, verified by a single kernel, crossing from
foundational set theory to program correctness. The tape is a set of ordered
pairs. The head position is a natural number constructed via the recursion
theorem. The correctness statement connects arithmetic ($\text{Plus}(a,b,c)$)
to computation ($\text{TMReaches}(\delta, c_0, n, c_f)$).

This depth is harder to achieve than breadth. Breadth (more theorems at the
same level) reuses established patterns. Depth (a new level on top of all
previous levels) introduces new proof engineering challenges at every layer,
and the proofs grow because they carry the entire pyramid beneath them.

The restriction to Z was a deliberate design choice. The LLM's initial
proof attempts used these stronger axioms freely; the human architect
imposed the restriction to Z, requiring the LLM to find alternative proofs.
This required developing non-trivial substitutes: successor injectivity was
reproved via transitive sets (Lemma~4.2.4 of the textbook) rather than
Regularity, and the recursion theorem's graph construction used power set
bounding ($\mathcal{P}(\mathcal{P}(\omega \cup \bigcup\bigcup f))$) rather
than Replacement. The result is a tighter axiom dependency than is typical
even in dedicated set theory formalizations --- all 203 theorems use only
Extensionality, Empty Set, Pairing, Union, Power Set, Separation, and
Infinity.

This illustrates the division of roles: the human sets the standard (use only
Z), the LLM finds proofs that meet it (alternative constructions), and the
kernel verifies the result (no Replacement, Regularity, or Choice in the
axiom dependencies).

% ============================================================
\section{Related Work}
\label{sec:related}
% ============================================================

\paragraph{LLM-driven theorem proving.}
AlphaProof~\citep{alphaproof2024} combines a language model with AlphaZero-style
reinforcement learning to solve competition mathematics in Lean. GPT-f~\citep{polu2020generative}
and follow-up work demonstrate LLM-based proof step generation in formal
systems. These approaches use LLMs within existing proof assistants (Lean, Isabelle)
and focus on search within a fixed logical framework. Our work differs in
building the logical framework itself, using ZFC rather than dependent type theory,
and targeting cross-abstraction reasoning rather than single-level proof search.

\paragraph{Interactive theorem provers.}
Lean/Mathlib~\citep{mathlib2020}, Coq~\citep{coq}, and Isabelle~\citep{isabelle}
are the dominant platforms for formalized mathematics. They are based on type theory
(CIC or simple type theory) and rely on human-driven interactive proof construction.
Mathlib contains over 100,000 theorems but has required a large community over
many years. Our approach substitutes LLM assistance for most human effort, at the
cost of a less mature library.

\paragraph{Set-theoretic formalization.}
Metamath's set.mm~\citep{metamath} is the largest ZFC-based formal library
($\sim$40,000 theorems). Metamath uses a minimal substitution-based proof
system; our use of Gentzen's sequent calculus (LK) provides a more natural
proof-theoretic framework with well-studied meta-properties (cut elimination,
the subformula property). Mizar~\citep{mizar} uses a set-theoretic foundation
with a richer type system. Both rely entirely on human proof authors. Our
contribution is demonstrating that LLM assistance can dramatically accelerate
ZFC formalization while maintaining full kernel verification.

\paragraph{LLM as proof search heuristic.}
The use of neural networks as heuristics in formal reasoning has precedent in
SAT solving~\citep{selsam2019learning} and program synthesis. Our architecture
makes this explicit: the LLM is an oracle that proposes proof steps, and the
kernel is the verifier that accepts or rejects them. The LLM's reliability is
irrelevant to the soundness of the final proof.

% ============================================================
\section{Discussion}
\label{sec:discussion}
% ============================================================

\paragraph{Limitations.}
This work demonstrates the pipeline on one textbook, one domain, and one
proof chain. Generality to other textbooks and domains remains to be shown.
The LLM was not fully autonomous --- a human architect made high-level design
decisions (theorem scoping, vocabulary structure, proof decomposition into
phases). The current theorem library, while spanning 10 abstraction levels,
covers only basic set theory and arithmetic.

\paragraph{The role of the textbook.}
Any scientific textbook that models its content mathematically is a potential
input to this pipeline. The textbook provides what the LLM lacks: strategic
direction, the right order of development, and conceptual organization. The
LLM provides what manual formalization lacks: the ability to work at the
scale of mechanical detail that formal proof requires. The ``easy to see''
steps in a textbook are exactly what LLMs are good at filling in --- and the
kernel catches the cases where ``easy to see'' is wrong.

\paragraph{Why ZFC.}
Physics, engineering, and applied mathematics express themselves in sets,
functions, and real numbers. These are native objects in ZFC. Encoding them
in dependent type theory is possible but introduces a translation layer that
adds complexity without expanding the class of problems that can be addressed.
For a system targeting the formalization of scientific knowledge broadly, ZFC
is the natural choice.

\paragraph{Depth as the hard direction.}
The demonstration in this paper prioritized depth (axioms to program
correctness) over breadth (many theorems at one level). We believe this is
the more important direction: depth demonstrates cross-abstraction reasoning,
which is the capability that existing tools lack. Breadth can be added
incrementally once the depth is established.

% ============================================================
\section{Conclusion}
\label{sec:conclusion}
% ============================================================

Formal proofs are too large for humans. This is not a temporary limitation of
current tools --- it is the nature of formal verification. The gap between a
two-line definition and its 11,500-line formal proof is inherent.

We have presented a pipeline that addresses this gap: textbooks provide
strategy, LLMs provide mechanical detail, and a minimal kernel provides
correctness guarantees. We demonstrated this pipeline by constructing a
verified proof chain from ZFC axioms to Turing machine correctness ---
203 theorems, 39,059 lines of proof, 10 abstraction levels, verified by
685 lines of trusted code.

The kernel caught a definition error in the guiding textbook that three years
of human review had missed. This is what formal verification is for: not
trusting humans, not trusting LLMs, but trusting a small, auditable kernel.

The formalized knowledge base --- the theorems, their dependencies, the
verified reasoning chains --- is a cumulative asset. Each theorem makes
future proofs easier. The library grows, and the LLM navigates it more
effectively as it grows. This is the beginning of a scalable approach to
formalizing the mathematical knowledge that underlies science and engineering.

\bibliographystyle{plainnat}
\bibliography{refs}

\end{document}
